\subsection{The source bar coalgebra and the compact comparison}
\label{subsec:koszul-source}

The hybrid correspondences must first define an algebra in a specified
coefficient category. Only then does its associative bar have a
meaning. We use the unitalized Ran category
$(\mathcal R_E^+,\star,K)$ of Section~\ref{sec:bkm-ran-envelope},
with right differential-module coefficients, cross-block chiral
convolution and empty unit $K$. A realization of constructible Hall
coefficients in this category is additional geometric data.
It must carry the products, orientation transport and higher collision
coherences, including wrapped and mixed configurations.

Write $V_E=\mathsf{Den}_{\Delta_5,E}$ for the BD current envelope
and $\mathcal A_{\Delta,E}=U_\star(\iota_*V_E)$ for its associative
Ran envelope. The associative comparison in this subsection has
$\mathcal A_{\Delta,E}$ as its target. A comparison directly with
$V_E$ as a BD chiral algebra is a different statement and needs the
corresponding chiral functor. The bar--cobar equivalence in
Proposition~\ref{prop:bkm-ran-bar-envelope} does not identify these
two recipients.

\begin{lemma}[Source and target bars]
\label{lem:source-target-koszul-separation}
For an augmented associative target $T$ with $T_\varnothing=\C$,
the counit $\Omega BT\to T$ is an internal equivalence in the
specified Ran category. It constructs neither a compact-source
algebra nor a morphism from such an algebra to $T$.
\end{lemma}
\begin{proof}
The counit depends functorially on the multiplication and
augmentation of $T$. Its construction has no input consisting of
compact correspondences, a realization of their coefficients,
orientations or a source-to-target morphism. In particular,
$C_X:=BT$ would define a coalgebra from $T$, without identifying it
with a coalgebra obtained from the compact geometry.
\end{proof}

\subsubsection*{Coefficient realization and the source coalgebra}

At a stage $R$, denote the realized hybrid algebra by $A_{X,R}$.
The index may bound geometric complexity or configuration arity.
It does not by itself specify a root-height quotient. Whenever a
target stage $T_R$ is used, it must be an actual augmented
associative object with supplied transition maps. A projection to
symmetric positive and negative root degrees is not such an object,
by Proposition~\ref{prop:bkm-positive-completion}. Identification
of the limit of the $T_R$ with $\mathcal A_{\Delta,E}$ is a
separate comparison.

\begin{definition}[Chiral Koszul source datum]
\label{def:chiral-koszul-source-datum}
For fixed stability, support and stage data $(\sigma,S,R)$, a
source datum $\mathfrak K^{\mathrm{Kos}}_{X,R}$ consists of the
following objects and morphisms in $\mathcal R_E^+$.
\begin{enumerate}
\item A realization of
$\mathcal F^{\mathrm{hyb}}_{X,\sigma,S,\leq R}$ as an augmented
associative algebra $A_{X,R}$ with empty component $\C$.
Its multiplication and all coherent collision operations come from
the supplied hybrid correspondences and orientations.
\item A coaugmented counital coassociative coalgebra $C_{X,R}$
with empty component $\C$, differential $d_C$, coproduct
$\Delta_R^{\mathrm{ch}}$ and an exhaustive filtration
$F^{\mathrm{co}}_{R,\bullet}$. The reduced coproduct is locally
nilpotent on each finite configuration and retained filtration
piece. If a global finite bound on that stage is used, it is included
in the datum.
\item A morphism of coalgebras
$b_{X,R}:C_{X,R}\to BA_{X,R}$ compatible with charge descent,
coaugmentation, counit and the supplied geometric operations.
Its cone is acyclic. The geometric construction of $C_{X,R}$
and this comparison are specified on local, wrapped and mixed
inputs, with all higher-colored collision coherences.
\item Compatible transition morphisms on these objects. Their
coherent diagrams include the reduced $E$-quotient and the
normal-ordered charge homomorphism $\overline\Pi_X$.
\end{enumerate}
A source-to-target comparison additionally supplies an augmented
associative morphism $\vartheta_R:A_{X,R}\to T_R$ and proves
its cone acyclic. Define
\[
 \varepsilon^{\mathrm{src}}_{\mathrm{bar},R}
 :=\varepsilon_{A_{X,R}}\circ\Omega b_{X,R}:
 \Omega C_{X,R}\longrightarrow A_{X,R},
 \qquad
 \Theta_{\mathrm{Kos},R}
 :=\vartheta_R\circ\varepsilon^{\mathrm{src}}_{\mathrm{bar},R}.
\]
The symbols $B,\Omega$ in this datum always refer to $\star$ and $K$.
\end{definition}

The finite binary flag atlas supplies particular associativity
and coproduct diagrams. It does not replace the higher-colored,
unit, refinement and quotient coherences required by this
definition. Nor does the projection-finite local sector provide the
coefficient realization of wrapped inputs. A restriction to ordered
words records the word-cut part of $C_{X,R}$. Comparison with the
full Ran coalgebra must preserve every collision operation.

The obstruction notation can retain these separate requirements:
\[
\begin{split}
 \mathfrak O_{\mathrm{Kos},R}=
 (&\mathfrak o^C_R,\mathfrak o^{\mathrm{fil}}_R,
 \mathfrak o^{\mathrm{conil}}_R,\mathfrak o^{\Delta,\mathrm{ch}}_R,
 \{\mathfrak o^{C,\mathrm{tr}}_{R'R}\},
 H^\bullet\operatorname{Cone}b_{X,R},\mathfrak o^\Omega_R,\\
 &\mathfrak o^{\mathrm{wrKos}}_R,\mathfrak o^{\mathrm{hyb}}_R,
 \mathfrak o^\Pi_R,\mathfrak o^{\mathrm{sep}}_R,
 \mathfrak o^{\Prim}_R).
\end{split}
\]
Here the first entries mean the stated coalgebra, filtration and
transition identities. $\mathfrak o^{\mathrm{wrKos}}_R$ includes
the wrapped coefficient realization and its bar domain.
The last entries require compatible charge descent, separation of
source from target, and any specified ordinary primitive realization.
An absent object or map is an unresolved construction, rather than
a numerical value of this tuple.

\begin{proposition}[Finite-length source bar]
\label{prop:canonical-source-bar-coalgebra}
Suppose the realized augmented algebra $A_{X,R}$ is supplied and
its positive augmentation ideal $a_{X,R}$ has no nonzero retained
$\star$-word of length greater than $N_R$. Then its normalized
bar coalgebra is
\[
 C^{\mathrm{bar}}_{X,R}=BA_{X,R}
 =K\oplus\bigoplus_{1\leq p\leq N_R}(sa_{X,R})^{\star p},
\]
with the internal and multiplication differentials. The suspension
$s$ is the bar suspension. Multiplication has degree one after
suspension. The coproduct is deconcatenation:
\[
 \Delta[x_1|\cdots|x_p]
 =1\otimes[x_1|\cdots|x_p]+[x_1|\cdots|x_p]\otimes1
  +\sum_{i=1}^{p-1}[x_1|\cdots|x_i]\otimes[x_{i+1}|\cdots|x_p].
\]
It is conilpotent and supplies a source-bar model with $b_{X,R}$
the identity.

Without a uniform $N_R$, the same construction is
finite at each configuration by the positive arity bound.
\end{proposition}
\begin{proof}
The normalized associative bar construction in this tensor category
has the displayed terms. Associativity with its coherent form gives
the square-zero bar differential. Cutting words gives a counital
coassociative coproduct compatible with it. Each reduced cut lowers
the available length, so more than $N_R$ nonempty blocks cannot
occur.

On a configuration $I$, positivity gives the bound $|I|$
instead. These are coefficient-category statements. To identify
the coproduct with geometric cutting correspondences requires the
realization and coherence maps stipulated in the source datum.
\end{proof}

If all nonempty geometric colors have positive HN height, the
bound on total HN height can imply $N_R\leq R$. That is a hypothesis
on the retained geometric algebra, not a conclusion about the full
two-sided root algebra. Additivity of $\overline\Pi_X$ ensures that
a bar word, each cut and each multiplication term have the same
descended Gram degree. This proves compatibility of gradings. It
constructs neither the realization nor the comparison $\vartheta_R$.

\begin{corollary}[The remaining source-to-target map]
\label{cor:source-bar-remaining-cobar-obstruction}
For a source datum in the specified Ran category,
$\varepsilon^{\mathrm{src}}_{\mathrm{bar},R}$ is an equivalence.
Hence $\Theta_{\mathrm{Kos},R}$ is an equivalence exactly when
$\vartheta_R$ is. A bar construction in another coefficient category
requires its own counit theorem before this conclusion applies.
\end{corollary}
\begin{proof}
The coalgebra map $b_{X,R}$ is an equivalence by assumption.
Apply the associative cobar and the bar--cobar counit of
Proposition~\ref{prop:bkm-ran-bar-envelope}. The two-out-of-three
property proves the assertion. A target-internal counit supplies
no part of $\vartheta_R$.
\end{proof}

For a strict representative, the comparison equations include
\[
 d_T\vartheta_R=\vartheta_Rd_X,\quad
 \vartheta_R\mu_X=\mu_T(\vartheta_R\star\vartheta_R),\quad
 \vartheta_R\eta_X=\eta_T,\quad
 \vartheta_R\rho^X_{R'R}=\rho^T_{R'R}\vartheta_{R'}.
\]
For coherent representatives, these equations are replaced by the
specified coherent homotopies. In either case the full coefficient
cone $\operatorname{Cone}\vartheta_R$ must be acyclic.
The Lie--commutative chiral theorem of
\cite[Theorems~5.1.1--5.2.1]{FrancisGaitsgoryCKD} does not supply
this compact associative comparison. It enters the envelope
construction at precisely the scope proved in
Proposition~\ref{prop:bkm-ran-bar-envelope}.

\subsubsection*{Compatibility with the geometric structures}

\begin{definition}[Finite Koszul comparison datum]
\label{def:finite-koszul-comparison-datum}
A comparison datum $\mathfrak K^{\mathrm{cmp}}_{X,R}$ supplements
the source datum by the following maps and commuting diagrams.
\begin{enumerate}
\item[\textup{(K0)}] The source counit and $\vartheta_R$ have
acyclic cones. If finite filtrations are used to prove this, the
associated-graded cone comparisons are acyclic at every retained
coefficient and parity.
\item[\textup{(K1)}] Every specified Weyl transport is intertwined
by $\Theta_{\mathrm{Kos},R}$. A reflection may map between chambers
or require a larger finite operation domain. Both domain and
codomain are specified. In chain representatives the equation is
$\Theta\tau_i^X-\tau_i^T\Theta=dH_i^W+H_i^Wd$.
The $H_i^W$ carry the higher coherences for their defined
compositions. No finite symmetric root set is assumed invariant
under the whole Weyl group.
\item[\textup{(K2)}] On the perfect complexes defining the geometric
Pfaffian and determinant lines, a supplied line isomorphism
intertwines orientation transport. Its type-II reflection sign is
$\epsilon_o(s_{\delta_i})=\widetilde\nu_{\Delta_5}(s_{\delta_i})=-1$.
The wall and chamber isomorphisms satisfy the cocycle identities of
$(O1)^+$ and $(O2_{\mathrm{atlas}})$.
\item[\textup{(K3)}] The associative map preserves products and
units, and its bar preserves coproducts, coaugmentations and
counits. If an ordinary Hall/Hopf realization is used to extract
primitives, it is separately supplied and its operation diagrams
commute. The invariant Lie form on those primitives is intertwined
with $B_{\mathfrak g}$. It is not identified with a pairing adjoint
to the ordinary primitive Hopf coproduct. Such an identification
is excluded by Proposition~\ref{prop:bkm-hopf-pairing-obstruction}.
\item[\textup{(K4)}] The primitive comparison carries the invariant source Lie radical
onto the target radical. It induces the comparison on their quotients.
When the target is $\mathfrak g$, its invariant form has zero
radical. Kernel and quotient squares are required in the actual
coefficient realization.
\item[\textup{(K5)}] This primitive comparison is the same map
$\iota_R^{\Prim}$ used in recognition. The induced maps of ordinary
enveloping algebras preserve PBW length and the finite weighted
pieces of Proposition~\ref{prop:bkm-finite-filtered-pbw}.
\item[\textup{(K6)}] Stage transitions and all the preceding
structures form coherent diagrams. In a strict chain model,
$\Theta_R\Omega(\rho^C_{R'R})-\rho^T_{R'R}\Theta_{R'}
=dH^{\mathrm{tr}}_{R'R}+H^{\mathrm{tr}}_{R'R}d$, with the
composition homotopies also specified.
\item[\textup{(K7)}] Completion uses the homotopy inverse limit
of these diagrams. To use ordinary inverse limits, require degreewise Mittag--Leffler
chain groups for the objects and comparison cones. Alternatively,
supply another explicit comparison between ordinary and homotopy
limits. Bar and cobar
commute with the chosen completed construction only when that
additional comparison has been proved.
\end{enumerate}
The notation $\mathfrak O_{\mathrm{Kcmp},R}$ denotes the two cone
cohomologies together with the Weyl, line, operation, radical, PBW
and transition defects in these diagrams. Their vanishing includes
the supplied coherent nullhomotopies, not just equal matrix ranks.
\end{definition}

\begin{proposition}[Compatible finite comparison]
\label{prop:finite-koszul-comparison-compatibility}
Under \textup{(K0)--(K6)}, the map
$\Theta_{\mathrm{Kos},R}:\Omega C_{X,R}\to T_R$ is an augmented
associative equivalence compatible with the specified Weyl,
orientation, primitive, invariant-form and enveloping diagrams.
\end{proposition}
\begin{proof}
Clause (K0) and the source counit give the equivalence.
Each remaining compatibility is a supplied commuting square or
coherent homotopy, and composition with the counit preserves it.
The radical comparison descends to quotient spaces. The ordinary
enveloping universal property gives the PBW comparisons.
This conclusion uses the stated realizations of these structures,
not an automatic extraction of a Hall object from the target.
\end{proof}

\begin{proposition}[Agreement on selected generators]
\label{prop:primitive-restriction-not-koszul}
Agreement with a prescribed map on a selected family of root
vectors does not imply an equivalence of associative coefficient
objects. For a supplied comparison $\vartheta_R$, the remaining
condition is acyclicity of its full coefficient cone.
\end{proposition}
\begin{proof}
The failure already occurs for ordinary augmented dg algebras.
Let $B$ be such an algebra and let $M$ be a dg $B$-bimodule with
nonzero cohomology. The square-zero extension $A=B\oplus M$ has
\[
 (b,m)(b',m')=(bb',bm'+mb'),\qquad M^2=0.
\]
Its augmented projection $p:A\to B$ fixes every selected vector
in the copy of $B$, but
$\operatorname{Cone}(p)\simeq M[1]$ is not acyclic.
This example asserts neither a connected Hopf structure on $A$
nor agreement on all its primitives. Under those additional
hypotheses, Proposition~\ref{prop:bkm-connected-hopf-recognition}
gives the positive ordinary Hopf recognition theorem.
\end{proof}

To establish the analogous conclusion in a chiral coefficient
category one must inspect the local comparison cone, not merely
the ordinary example. Additional wrapped, mixed or descendant
coefficients can lie outside the selected root-vector family.
The invariant form and the ordinary primitive map likewise need
the specified realizations in (K3)--(K5).

\begin{proposition}[Passage to completion]
\label{prop:chiral-koszul-source-datum-consequence}
Suppose source and comparison data are supplied on a cofinal
sequence of stages, with the coherent transitions of (K6).
Then their homotopy inverse limits give equivalences
\[
 \operatorname*{holim}_R C_{X,R}
       \simeq\operatorname*{holim}_R BA_{X,R},\qquad
 \operatorname*{holim}_R\Omega C_{X,R}
       \simeq\operatorname*{holim}_R T_R,
\]
with all specified compatible operations.
Under (K7), the same assertions hold for the ordinary inverse
limits of the chosen chain models. Writing these limits as
$B A_X$, $\Omega C_X$ and $\mathcal A_{\Delta,E}$ additionally
requires the corresponding bar, cobar and target-limit comparisons.
\end{proposition}
\begin{proof}
A homotopy limit is a functor on the coherent diagram category and
preserves equivalences. Apply it to $b_{X,R}$ and
$\Theta_{\mathrm{Kos},R}$ and their compatibility diagrams.
The ordinary-limit assertion follows from
Proposition~\ref{prop:bkm-limit-complexes}.
For actual degreewise finite-dimensional towers, the required
Mittag--Leffler property is automatic once their transition maps
exist. In general it is not automatic, as the polynomial-tail
example there shows. Nothing in preservation of equivalences by
limits alone says that bar or cobar commutes with that limit.
\end{proof}

\subsubsection*{Twisting cochains and the boundary algebra}

A twisting cochain from $C_{X,R}$ to an augmented algebra has
unsuspended degree one. It satisfies
\[
 d\tau+\tau d+\mu(\tau\star\tau)\overline\Delta=0,
\]
with the suspension signs fixed by the bar convention.
Equivalently it is a morphism from the cobar. For a bar source it
includes the canonical twisting cochain. This equation describes
coupling data at the specified carrier and tensor. It does not
identify the bar coalgebra with the bulk algebra.

A proposed bulk center belongs to the associative boundary algebra
$A_b=R\operatorname{End}(b)$ and requires its separate map and
hypotheses, as in Remark~\ref{rem:bkm-universal-holography-scope}.
The empty coefficient $K$, the BD vacuum, the augmentation of
$A_{X,R}$ and a prepared state are not interchangeable. In
particular, obtaining a completed coefficient equivalence to
$\mathcal A_{\Delta,E}$ still leaves the state comparison and its
protected trace to be constructed. The geometric primitive
recognition in the next subsection is another operation-preserving
comparison, with its own explicit source data.
